% Generated draft: requires mathematical and source review.
\documentclass[11pt]{article}
\usepackage[margin=25mm]{geometry}
\usepackage{fontspec}
\setmainfont{Latin Modern Roman}
\usepackage{amsmath,amssymb,mathtools,unicode-math}
\setmathfont{Latin Modern Math}
\usepackage{enumitem,needspace,xurl,hyperref}
\setlength{\parindent}{0pt}
\setlength{\parskip}{5pt}
\setlength{\emergencystretch}{3em}
\newcommand{\N}{\mathbb{N}}
\newcommand{\Z}{\mathbb{Z}}
\newcommand{\Q}{\mathbb{Q}}
\newcommand{\R}{\mathbb{R}}
\newcommand{\C}{\mathbb{C}}
\newcommand{\F}{\mathbb{F}}
\newcommand{\A}{\mathbb{A}}
\newcommand{\PP}{\mathbb{P}}
\newcommand{\E}{\mathbb{E}}
\newcommand{\eps}{\varepsilon}
\newcommand{\dd}{\,\mathrm{d}}
\begin{document}
\section*{UnsolvedMath ID: 701320}
\textbf{Source problem code:} catalogue entry
\subsection{Definitions and mathematical statement}
Interpret compact groups as compact Hausdorff topological groups, and embeddings as continuous group homomorphisms that are homeomorphisms onto their images. The supplied source does not explicitly state the Hausdorff convention.

Let \(\mathbf{A}\) be the class of compact groups \(A\) with the following amalgamation property: for every pair of compact groups \(B,C\) and every pair of embeddings
\[
i_B:A\hookrightarrow B,\qquad i_C:A\hookrightarrow C,
\]
there exist a compact group \(D\) and embeddings
\[
j_B:B\hookrightarrow D,\qquad j_C:C\hookrightarrow D
\]
such that
\[
j_B\circ i_B=j_C\circ i_C.
\]
No requirement is imposed that \(j_B(B)\cap j_C(C)\) equal the common image of \(A\).

Let \(\mathbb{T}=\R/\Z\), with its additive group structure and quotient topology. For a topological group \(A\), let \(A^{\circ}\) denote the connected component of its identity, equipped with the subspace topology. The three questions of G. M. Bergman are:
\begin{enumerate}
\item[(a)] Is \(\mathbb{T}\in\mathbf{A}\)?
\item[(b)] Does there exist a nonabelian connected compact Lie group \(G\) such that \(G\in\mathbf{A}\)?
\item[(c)] Is it true that
\[
\forall A\in\mathbf{A},\qquad A^{\circ}\in\mathbf{A}?
\]
\end{enumerate}

As background, the supplied statement attributes to Bergman's 1987 paper the results that every member of \(\mathbf{A}\) is a finite-dimensional compact Lie group, not necessarily connected, satisfying a strong ``local simplicity'' property, and that every finite group belongs to \(\mathbf{A}\). The source does not define that additional property; it is retained here as attributed background and is not an extra hypothesis in the definition of \(\mathbf{A}\).
\subsection{Short English statement}
The class \(\mathbf{A}\) consists of compact groups over which any two compact group extensions can be embedded compatibly into one compact group. Bergman asks whether the circle group belongs to this class, whether any nonabelian connected compact Lie group belongs to it, and whether membership passes to the identity component.
\subsection{Sources}
\begin{itemize}
\item Local UnsolvedMath registry: permanent ID \(701320\), exact source code \texttt{catalogue entry}, title \emph{Kourovka Notebook Problem 17.18}. The supplied record labels the problem \emph{open}; this label is supplied registry information, not an independently verified assessment.
\item \emph{The Kourovka Notebook (Unsolved Problems in Group Theory)}, Problem 17.18, proposed by G. M. Bergman, 17th issue (2010). The supplied background also identifies the 21st issue edition (2026), \texttt{arXiv:1401.0300v48}. The collection description names E. I. Khukhro and V. D. Mazurov as editors.
\item G. M. Bergman, paper cited in the supplied statement as \emph{Manuscr. Math.}, volume \(58\) (1987), pp.\ 253--281. No article title is supplied. The stated background results are attributed to this reference by the source and have not been independently verified here.
\end{itemize}
\end{document}
